%%% Please fill in basic information on your thesis, which will be automatically
%%% inserted at the right places. You need to replace \xxx{...} by real data.

% Type of your thesis:
%	"bc" for Bachelor's
%	"mgr" for Master's
%	"phd" for PhD
%	"rig" for rigorosum
\def\ThesisType{bc}

% Language of your study programme:
%	"cs" for Czech
%	"en" for English
\def\StudyLanguage{en}

% Thesis title in English (exactly as in the official assignment)
% (Note: \xxx is a "ToDo label" which makes the unfilled visible. Remove it.)
\def\ThesisTitle{Procedural Terrain Generation with Fixed Points}

% Author of the thesis (you)
\def\ThesisAuthor{Oliver Wakeford}

% Year when the thesis is submitted
\def\YearSubmitted{2025}

% Name of the department or institute, where the work was officially assigned
% (according to the Organizational Structure of MFF UK in English,
% see https://www.mff.cuni.cz/en/faculty/organizational-structure,
% or a full name of a department outside MFF)
\def\Department{Software and Computer Science Education}

% Is it a department (katedra), or an institute (ústav)?
\def\DeptType{department}

% Thesis supervisor: name, surname and titles
\def\Supervisor{Petr Šimůnek}

% Supervisor's department (again according to Organizational structure of MFF)
\def\SupervisorsDepartment{Software and Computer Science Education}

% Study programme (does not apply to rigorosum theses)
\def\StudyProgramme{BSc. Computer Science}

% An optional dedication: you can thank whomever you wish (your supervisor,
% consultant, who provided you with tea and pizza, etc.)
\def\Dedication{
I want to express my gratitude to my supervisor for his guidance, encouragement, and support throughout the development of this thesis. I am also sincerely thankful to my friends and family, whose belief in me and constant support have made this journey not only possible, but meaningful and enjoyable.
}

% Abstract (recommended length around 80-200 words; this is not a copy of your thesis assignment!)
\def\Abstract{
This thesis tackles the challenge of procedural terrain generation with fixed height constraints in Unity. We developed a system that creates natural-looking terrain around areas that must remain unchanged, identified through a mask image. Unlike other approaches, our method first generates terrain in modifiable regions, then applies specialised post-processing to create gentle integration with the fixed points. The core innovation is a distance-based algorithm that adjusts terrain features based on their proximity to fixed points, promoting seamless transitions. The system is implemented as a Unity component, enabling users to apply and customise terrain operations directly on terrain objects. To streamline parameter tuning, we introduce an Interactive Genetic Algorithm that allows users to evolve terrains through visual selection rather than manual input. The resulting landscapes exhibit cohesive transitions between fixed and generated areas, creating realistic formations that could conceivably exist in the natural world.
}

% 3 to 5 keywords (recommended) separated by \sep
% Keywords are useful for indexing and searching for the theses by topic.
\def\ThesisKeywords{%
procedural terrain generation\sep fixed height constraints\sep Unity}

% If any of your metadata strings contains TeX macros, you need to provide
% a plain-text version for use in XMP metadata embedded in the output PDF file.
% If you are not sure, check the generated thesis.xmpdata file.
\def\ThesisAuthorXMP{\ThesisAuthor}
\def\ThesisTitleXMP{\ThesisTitle}
\def\ThesisKeywordsXMP{\ThesisKeywords}
\def\AbstractXMP{\Abstract}

% If your abstracts are long and do not fit in the infopage, you can make the
% fonts a bit smaller by this setting. (Also, you should try to compress your abstract more.)
\def\InfoPageFont{}
%\def\InfoPageFont{\small}  % uncomment to decrease font size

% If you are studing in a Czech programme, you also need to provide metadata in Czech:
% (in English programmes, this is not used anywhere)

\def\ThesisTitleCS{\xxx{Název práce česky}}
\def\DepartmentCS{\xxx{Název katedry česky}}
\def\DeptTypeCS{\xxx{Katedra}}
\def\SupervisorsDepartmentCS{\xxx{katedra vedoucího}}
\def\StudyProgrammeCS{\xxx{studijní program}}

\def\ThesisKeywordsCS{%
\xxx{klíčová slova\sep klíčové fráze}
}

\def\AbstractCS{%
\xxx{Abstrakt práce přeložte také do češtiny.}
}
